% !Mode:: "TeX:UTF-8"
% 任务书中的信息
%% 原始资料及设计要求
\assignReq
{本设计的原始资料是PeMS04、PeMS07、PeMS08等公开的交通流量数据集}
{数据集转换成适配LibCity框架的原子文件格式，包含流量、占有率等数据}
{该设计采用慢思考—快反应的双层架构，支持自然语言交互和结构化输出}
{将大语言模型和时空预测模型协同，使非专业人员也可通过自然语言操作}
{最终实现"自然语言—工具调度—数据分析—报告生成"闭环。}

%% 工作内容
\assignWork
{本设计实现了一个基于大语言模型的交通智能体项目以本地千问推理引擎}
{为核心，通过LangGraph框架构建ReAct推理循环集成流量查询、分析和}
{预测三个工具，根据用户指令先进行意图识别再调用工具返回结果，数据}
{层面基于PeMS08数据集构建数据库，通过FastAPI提供后端微服务接}
{口，保障多工具协同的正确性。}

%% 参考文献
\assignRef
{[1]Jiang J, Han C, Zhao W X, et al. PDFormer: Propagation Delay-Aware }
{Dynamic Long-Range Transformer for Traffic Flow Prediction[C]. AAAI 2023}
{[2]李健, 朱国军, 王奥, 等. 城市交通知识增强大语言模型构建及应用探索[J].}
{城市交通, 2025, 23(02): 1-12+38}
{[3]奚志恒, 陈文翔, 郭昕, 等. 基于大语言模型的智能体：发展与未来展望[J].}
{中国科学: 信息科学, 2026, 56(2): 485-486}
{[4]李炎英, 王新宇, 王晓, 等. 基于大语言模型的交通异常事件检测与辅助决}
{策[J].智能科学与技术学报, 2024, 6(3): 347-355}
